\documentclass[fontsize=10.5pt]{jlreq}
\usepackage{mystyle}

\title{Deep Learning from Scratch Mathematical Questions generated by AI}
\author{}
\date{\today}

\begin{document}

\begin{abstract}
  これは「ゼロから作るDeep Learning」をはじめとする機械学習や深層学習において、理解を促進させるためにAIが設定した問題を解くためのものである。
\end{abstract}

\section{強化学習SARSA 重点サンプリングにおける疑問}

\begin{question}[TDターゲットの期待値の定義]

\end{question}
  TD学習の更新は、本質的に「TDターゲットの期待値」に向かって現在の価値を近づけていく操作です。

  ターゲット方策$\pi$にしたがう場合の理想的なTDターゲットの条件付き期待値$\mathbb{E}_\pi\left(R_t + rQ_\pi(S_{t+1}, A_{t+1}) \mid S_t = s, A_t = a\right)$を、環境の遷移確率$p(s' \mid s, a)$とターゲット方策$\pi(a' \mid s')$を用いて、和の記号$\sum$で書き下すとどのようになりますか？
\end{document}

\newpage

\begin{question}[挙動方策$b$への変換（重点サンプリングの導出）]
  問1で書き下した期待値の式において、サンプリングの元となる行動の分布を$\pi$から挙動方策$b$に変更したいと思います。期待値の数式を変形して無理やり$b(a' \mid s')$を出現させた場合、数式として帳尻を合わせるために、どのような分数（重み）が内部に掛けられることになりますか？
\end{question}

\newpage

\begin{question}[ご提示いただいた式の検証]
  問2までの数理的な考察を踏まえて、ご自身が最初に提示された以下の2つの要素を見直してみてください。
  \begin{enumerate}
    \item \textbf{重み$\rho$の添え字（時刻）：}$\displaystyle \rho = \frac{\pi(A_t \mid S_t)}{b(A_t \mid S_t)}$とされていますが、既に$A_t$が「条件」として実行済みである$Q(S_t, A_t)$の更新において、真に分数の補正が必要になるのはどの時刻の行動でしょうか？
    \item \textbf{$\rho$が乗算される範囲：}ご提示の式では$\rho \dot \left(R_t + \gamma Q_\pi(S_{t+1}, A_{t+1})\right)$となっていますが、報酬$R_t$は行動のサンプリング分布の違いによる補正の対象になるでしょうか？
  \end{enumerate}
\end{question}